\documentclass[10pt,a4paper]{amsart}
\usepackage[margin=19mm]{geometry}
\usepackage{amsmath,amssymb,mathtools}
\usepackage[scr=rsfs,cal=euler]{mathalfa}
\usepackage{xfrac}
\usepackage[unicode,hidelinks,bookmarks=false]{hyperref}
\newcommand{\ie}{i.e.\ }
\newcommand{\cf}{cf.\ }
\newcommand{\resp}{respectively\ }
\newcommand{\Subsection}{\subsection}
\providecommand{\nobreakdash}{\nobreak\hbox{-}}
\setlength{\emergencystretch}{2em}
\multlinegap0pt
\usepackage{mathtools}
\usepackage{amssymb}
\usepackage{stackengine}
\newtheorem{lemma}{Lemma}
\title{Special Dirichlet Processes: Structure, Uniqueness and Stability}
\author{Philip Kennerberg}
\date{Source: arXiv:2606.02056v1 (CC BY 4.0)}
\begin{document}
\maketitle

\noindent\emph{Source and licence.} Special Dirichlet Processes: Structure, Uniqueness and Stability. Version arXiv:2606.02056v1 (CC BY 4.0), distributed by arXiv under the Creative Commons Attribution 4.0 International licence, http://creativecommons.org/licenses/by/4.0/, as recorded in the arXiv licence metadata for this exact version and shown on the abstract page https://arxiv.org/abs/2606.02056v1. Full licence notice: \texttt{licenses/arxiv-2606.02056-CC-BY-4.0.txt}.\par\smallskip

\noindent\textit{Editorial scope.} Counted unit: the article's lemma lem:SD-decomposition-stability (source0.tex lines 651--672). The article's complete original proof of it is delivered with it (source0.tex lines 674--856, one \texttt{proof} environment of 183 lines). No mathematics is written, repaired or re-derived: the statement and the proof are the article's own lines, in the article's own notation, and no retained line is substituted or re-flowed.  No further source-local context is needed: every cross-reference of every retained line resolves inside the two delivered spans.  Nothing was omitted from any retained span, so the package carries no omission record.  No retained line contains a citation, so the wrapper carries no bibliography and no bibliography entry.  No retained line contains a figure environment, an image-inclusion command or drawing code, so no image is delivered and the package declares no asset dependency.  Editorial formatting only, in addition: an AMS \texttt{amsart} preamble that restates the article's own declarations and macros used by the retained lines, namely the environments \texttt{lemma} on separate counters, together with the standard packages those lines need.  The retained environments print in this document as Lemma 1 in order of appearance, whereas the article prints the same items with the numbering of its own full document.  The complete native source of the article as distributed in the arXiv e-print for this exact version is bundled unchanged as \texttt{sources/201873/source0.tex}.
\begin{lemma}[Stability of the Special--Dirichlet decomposition]
\label{lem:SD-decomposition-stability}
Let \(X^n\) and \(X\) be Special--Dirichlet processes on \([0,t]\), with
canonical decompositions
\[
        X^n=M^n+\Gamma^n,
        \qquad
        X=M+\Gamma,
\]
where \(M^n_0=M_0=0\). Assume that
\[
        [X^n-X]_t \xrightarrow{\mathbb P} 0 .
\]
Then
\[
        [\Gamma^n-\Gamma]_t \xrightarrow{\mathbb P} 0
\]
and
\[
        [M^n-M]_t \xrightarrow{\mathbb P} 0 .
\]
\end{lemma}

\begin{proof}
Set
\[
        Y^n:=X^n-X,\qquad
        N^n:=M^n-M,\qquad
        B^n:=\Gamma^n-\Gamma .
\]
Then
\[
        Y^n=N^n+B^n .
\]
By the definition of Special--Dirichlet processes, \(B^n\) has zero
continuous quadratic variation and its jump times are carried by a thin
predictable set. Moreover,
\[
        \Delta B^n_s\in\mathcal F_{s-}
\]
at every jump time of \(B^n\).

Let \((T^n_k)_{k\ge1}\) be predictable stopping times exhausting the jump
times of \(B^n\). Since
\[
        [B^n]^c_t=0,
\]
we have
\[
        [B^n]_t
        =
        \sum_{k\ge1:T^n_k\le t}
        \bigl(\Delta B^n_{T^n_k}\bigr)^2 .
\]

Fix \(\eta>0\), and define the stopping time
\[
        \tau^n_\eta
        :=
        \inf\{s\le t:[Y^n]_s>\eta\}.
\]
On the event \(\{\tau^n_\eta>t\}\), we have \([Y^n]_t\le \eta\).

We first estimate the predictable jump part before \(\tau^n_\eta\). After
standard localization, we may assume that the local martingale \(N^n\) is
square-integrable up to \(t\). For every \(k\), on the event
\(\{T^n_k<\tau^n_\eta\}\), the variable
\[
        \Delta B^n_{T^n_k}
\]
is \(\mathcal F_{T^n_k-}\)-measurable. Since \(N^n\) is a local martingale,
\[
        \mathbb E\!\left[
        \Delta N^n_{T^n_k}
        \mid
        \mathcal F_{T^n_k-}
        \right]
        =0 .
\]
Hence
\[
\begin{aligned}
\mathbb E\!\left[
        \bigl(\Delta Y^n_{T^n_k}\bigr)^2
        \mid
        \mathcal F_{T^n_k-}
        \right]
&=
\mathbb E\!\left[
        \bigl(\Delta N^n_{T^n_k}
        +
        \Delta B^n_{T^n_k}\bigr)^2
        \mid
        \mathcal F_{T^n_k-}
        \right]
\\
&=
        \bigl(\Delta B^n_{T^n_k}\bigr)^2
        +
        \mathbb E\!\left[
        \bigl(\Delta N^n_{T^n_k}\bigr)^2
        \mid
        \mathcal F_{T^n_k-}
        \right]
\\
&\ge
        \bigl(\Delta B^n_{T^n_k}\bigr)^2 .
\end{aligned}
\]
Therefore,
\[
        \mathbb E\!\left[
        \bigl(\Delta B^n_{T^n_k}\bigr)^2
        \mathbf 1_{\{T^n_k<\tau^n_\eta\}}
        \right]
        \le
        \mathbb E\!\left[
        \bigl(\Delta Y^n_{T^n_k}\bigr)^2
        \mathbf 1_{\{T^n_k<\tau^n_\eta\}}
        \right].
\]
Summing over \(k\) and using monotone convergence gives
\[
\begin{aligned}
\mathbb E\!\left[
        \sum_{k\ge1:T^n_k<\tau^n_\eta}
        \bigl(\Delta B^n_{T^n_k}\bigr)^2
        \right]
&\le
\mathbb E\!\left[
        \sum_{k\ge1:T^n_k<\tau^n_\eta}
        \bigl(\Delta Y^n_{T^n_k}\bigr)^2
        \right]
\\
&\le
\mathbb E\!\left[
        \sum_{s<\tau^n_\eta}
        \bigl(\Delta Y^n_s\bigr)^2
        \right]
\\
&\le
\eta .
\end{aligned}
\]
The last inequality follows from the definition of \(\tau^n_\eta\).

By Markov's inequality,
\[
\mathbb P\!\left(
        \sum_{k\ge1:T^n_k<\tau^n_\eta}
        \bigl(\Delta B^n_{T^n_k}\bigr)^2
        >\varepsilon
        \right)
        \le
        \frac{\eta}{\varepsilon}.
\]
Hence
\[
\begin{aligned}
\mathbb P\!\left([B^n]_t>\varepsilon\right)
&\le
\mathbb P(\tau^n_\eta\le t)
+
\mathbb P\!\left(
        \sum_{k\ge1:T^n_k<\tau^n_\eta}
        \bigl(\Delta B^n_{T^n_k}\bigr)^2
        >\varepsilon
        \right)
\\
&\le
\mathbb P([Y^n]_t>\eta)
+
\frac{\eta}{\varepsilon}.
\end{aligned}
\]
Since \([Y^n]_t=[X^n-X]_t\to0\) in probability, we obtain
\[
        \limsup_{n\to\infty}
        \mathbb P\!\left([B^n]_t>\varepsilon\right)
        \le
        \frac{\eta}{\varepsilon}.
\]
Letting \(\eta\downarrow0\), we conclude that
\[
        [B^n]_t=[\Gamma^n-\Gamma]_t
        \xrightarrow{\mathbb P}0 .
\]

It remains to prove the convergence of the martingale parts. Since
\[
        N^n=Y^n-B^n,
\]
we have, by the triangle inequality for quadratic variation,
\[
        [N^n]_t^{1/2}
        \le
        [Y^n]_t^{1/2}
        +
        [B^n]_t^{1/2}.
\]
Therefore
\[
        [M^n-M]_t=[N^n]_t
        \xrightarrow{\mathbb P}0 .
\]
\end{proof}

\end{document}
